% Programa del periodo para estudiantes · Filosofía 10.° · Periodo IV · PILOTO (borrador)
% Se entrega al inicio del periodo. Sale de malla-P4.tex (aprendizajes D1–D5) y de
% parcelacion-P4.tex (cronograma de 16 clases), en lenguaje para estudiantes.
% Compilar: bash curriculo/plantilla/compilar.sh curriculo/filosofia/10/programa-P4.tex
\documentclass[programa]{planeacion}
\materia{filosofia}
\grado{10.°} \periodo{IV} \anio{2027}
\clasesporsemana{2}
\carpeta{}   % enlace a la carpeta de Drive de esta planeación (se completa al compartirla)
\begin{document}
\fichaprograma

\destacado{PREGUNTA DEL\newline PERIODO}{¿Pueden convivir la fe y la razón? ¿Qué cambió en la
  manera de pensar al ser humano cuando Europa pasó de la Edad Media al Renacimiento?}

\seccion{1}{Lo que vas a aprender}
Al terminar el periodo podrás:
\begin{aprendizajes}
  \aprendizaje{D1}{Explicar cómo las ideas de Platón y Aristóteles llegaron a la Edad Media y
    ubicar a los grandes pensadores en una línea de tiempo.}
  \aprendizaje{D2}{Interpretar fragmentos de San Agustín y explicar cómo entendía la relación entre
    la fe y la razón.}
  \aprendizaje{D3}{Reconocer el método de la \emph{quaestio} de Tomás de Aquino y usarlo para
    discutir con rigor un problema actual.}
  \aprendizaje{D4}{Evaluar con evidencias si la Edad Media fue una «época oscura» y explicar qué
    trajeron la imprenta y el humanismo de Erasmo.}
  \aprendizaje{D5}{Interpretar la crítica social de la \emph{Utopía} de Tomás Moro y proponer y
    defender una utopía para tu comunidad.}
\end{aprendizajes}

\seccion{2}{Cómo vamos a trabajar}
En esta clase \textbf{no se dicta}: el docente explica y \textbf{tú tomas tus propios apuntes} en el
cuaderno; al final de esas clases se \textbf{firman}. \textbf{Cada clase tiene un entregable corto} en
el cuaderno (un mapa, una carta, una ficha, un veredicto…) que recibe un \textbf{sello}. Habrá
actividades distintas a lo habitual: una disputa medieval, un juicio a la Edad Media, una carrera
entre un copista y la imprenta. Lo importante es pensar, argumentar y participar.
\begin{itemize}\setlength{\itemsep}{1pt}
  \item \textbf{Trae siempre} tu cuaderno de filosofía: ahí quedan tus apuntes, sellos y firmas.
  \item \textbf{En los debates y el juicio} se discuten ideas, no personas; se pide la palabra y se escucha.
  \item \textbf{Si faltas}, ponte al día con los apuntes de un compañero y presenta el entregable la
    clase siguiente.
\end{itemize}

\seccion{3}{Cronograma de clases}
\begin{cronograma}
  \sesion{La herencia griega: \emph{La escuela de Atenas}.}{«¿Quién es quién?» en el cuadro de
    Rafael.}{\firma{Apuntes.}}
  \sesion{El viaje de los libros.}{Reto: traducir una frase de Aristóteles por tres lenguas.}{\sello{Mapa del viaje.}}
  \sesion{San Agustín: «vuelve a ti mismo».}{Explicación y lectura de fragmentos.}{\firma{Apuntes.}}
  \sesion{Una carta a Agustín.}{Le escribes desde 2027.}{\sello{Carta.}}
  \sesion{Fe y razón.}{Experimento de la caja cerrada.}{\sello{Reflexión.}}
  \sesion{Tomás de Aquino y la universidad.}{Cómo era estudiar en París en 1260.}{\firma{Apuntes.}}
  \sesion{Disputa medieval en vivo.}{Recreamos una \emph{disputatio}.}{\sello{Esquema de la disputa.}}
  \sesion[t]{Taller 1: mi \emph{quaestio}.}{En parejas, sobre un problema actual.}{Taller.}
  \sesion{Juicio a la Edad Media (preparación).}{Roles: fiscalía, defensa, jurado y testigos.}{\sello{Ficha del personaje.}}
  \sesion{Juicio a la Edad Media.}{Audiencia y veredicto.}{\sello{Veredicto razonado.}}
  \sesion{La imprenta.}{Carrera copista contra imprenta.}{\sello{Cálculo y conclusiones.}}
  \sesion[q]{Quiz del periodo.}{Agustín, Tomás y la Edad Media.}{Quiz.}
  \sesion{Erasmo: la Locura toma la palabra.}{Lectura dramatizada.}{\sello{Discurso de la Locura.}}
  \sesion{Tomás Moro y la \emph{Utopía}.}{Explicación y fragmentos.}{\firma{Apuntes.}}
  \sesion[t]{Taller 2: diseña tu isla de Utopía.}{Mapa, tres leyes y sus riesgos.}{Taller.}
  \sesion[e]{Evaluación de periodo.}{Prueba escrita sobre todo el periodo.}{Evaluación.}
\end{cronograma}

\seccion{4}{Evaluación}
Se evalúa el proceso durante todo el periodo, no solo la prueba final. Cada actividad se valora con
la escala institucional, equivalente a la nacional: \textbf{Superior}, \textbf{Alto},
\textbf{Básico} y \textbf{Bajo}.\par\vspace{4pt}
{\rowcolors{2}{white}{cremafila}\estilotabla\renewcommand{\arraystretch}{1.12}\small
\renewcommand{\tabularxcolumn}[1]{m{#1}}% contenido centrado en vertical
\begin{tabularx}{\linewidth}{|>{\RaggedRight\arraybackslash}X|>{\centering\arraybackslash}m{1.25cm}|>{\centering\arraybackslash}m{2.15cm}|>{\RaggedRight\arraybackslash}m{5.95cm}|}
  \hline
  \cabecera\textbf{\color{pinotexto}Actividad} & \textbf{\color{pinotexto}Semana} & \textbf{\color{pinotexto}Aprendizajes} & \textbf{\color{pinotexto}Qué se tiene en cuenta}\\ \hline
  Entregables de clase (sellos) y apuntes firmados & 1--7 & D1--D5 & Uno por clase, completo y entregado ese día.\\ \hline
  Taller 1: mi \emph{quaestio} & 4 & D3 & Estructura completa, objeciones serias, respuesta argumentada.\\ \hline
  Quiz del periodo & 6 & D1--D4 & Comprensión de lo visto hasta la semana 6.\\ \hline
  Taller 2: diseña tu isla de Utopía & 8 & D5 & Leyes con razones, riesgos previstos y exposición.\\ \hline
  Evaluación de periodo & 8 & D1--D5 & Análisis de textos, imágenes y argumentos.\\ \hline
  Autoevaluación & 8 & --- & Reflexión honesta sobre tu propio proceso.\\ \hline
\end{tabularx}\arrayrulecolor{black}}

\seccion{5}{Lecturas y recursos}
{\small Rafael, \emph{La escuela de Atenas} (1509--1511) · San Agustín, \emph{Confesiones}, libro I,
cap. 1 · Tomás de Aquino, \emph{Suma teológica} (un artículo) · Erasmo de Rotterdam, \emph{Elogio de la
locura} (1511) · Tomás Moro, \emph{Utopía} (1516), libros I y II. Los fragmentos se entregan en clase y
quedan, con las presentaciones y guías, en la carpeta del periodo (código QR).\par}

\recibido
\end{document}
